% Formalization of the optimization problem in the cloud server optimization project
\documentclass[11pt]{article}
\usepackage[utf8]{inputenc}
\usepackage[T2A]{fontenc}
\usepackage[russian,english]{babel}
\usepackage{amsmath,amssymb}
\usepackage{hyperref}

\title{Формализация задачи оптимизации\\(Cloud Server Optimization)}
\author{cloud\_server\_optimization}
\date{\today}

\begin{document}
\maketitle

\section*{1. Обозначения и входные данные}

\paragraph{Поды (реплики).} Пусть
\(\mathcal{P}=\{1,2,\dots,P\}\) --- множество реплик.
Каждая реплика \(p\in\mathcal{P}\) имеет:
\begin{align*}
  c_p &>0 &\text{(CPU-запрос)},\\
  m_p &>0 &\text{(память)}.
\end{align*}

\paragraph{Типы узлов (шаблоны).} Пусть
\(\mathcal{T}=\{1,2,\dots,T\}\) --- множество типов узлов.
Каждый тип \(t\in\mathcal{T}\) имеет параметры:
\begin{align*}
  C_t &>0 &\text{(CPU-ёмкость)},\\
  M_t &>0 &\text{(память)},\\
  P_t &\in\mathbb{Z}_{>0} &\text{(максимум подов)}.
\end{align*}

В реализации эти величины задаются через структуру \texttt{NodeTemplate}.

\section*{2. Переменные}

Для задания размещения вводятся бинарные переменные.

\paragraph{Размещение.} Для каждого пода \(p\) и каждого (конкретного) узла \(n\) вводим:
\[ x_{p,n}\in\{0,1\} \qquad(=1\text{, если }p\text{ запущен на }n). \]

\paragraph{Активность узла.} Для каждого узла \(n\) вводим:
\[ y_n\in\{0,1\} \qquad(=1\text{, если узел }n\text{ активен}). \]

\section*{3. Ограничения}

\paragraph{Каждый под размещён ровно в одном месте:}
\[
  \forall p:\quad \sum_n x_{p,n} = 1.
\]

\paragraph{Ограничения ресурсов на узле:}
\begin{align*}
  \forall n:\quad &\sum_p c_p\,x_{p,n} \le C_n, \\
  \forall n:\quad &\sum_p m_p\,x_{p,n} \le M_n, \\
  \forall n:\quad &\sum_p x_{p,n} \le P_n.
\end{align*}

\paragraph{Связь активности узла с размещением:}
\[
  \forall p,n:\quad x_{p,n} \le y_n.
\]

(Если на узле есть хотя бы один под, то он считается активным.)

\section*{4. Целевая функция}

Основная оптимизационная цель в проекте --- минимизировать число активных узлов:
\[
  \min \sum_n y_n.
\]

В моделях с учётом задержки (SLO) можно вводить дополнительные штрафы
(например, штраф за среднюю нагрузку выше порога или за превышение порога задержки).
В коде задействована нелинейная модель задержки:
\[
  \text{latency}(u) = 90 + 35u + 240\max(0, u - 0.75)^2,
\]
где \(u\) --- средняя загрузка CPU.

\section*{5. Связь с эвристикой в коде}

\begin{itemize}
  \item \textbf{FFD (First Fit Decreasing)} реализует жадную упаковку pod-ов по узлам.
  \item \textbf{Ребалансировка} (миграции) ограничена бюджетом миграций.
  \item \textbf{Прогнозирование (EWMA)} используется для расчёта требуемого количества активных узлов с запасом (headroom).
\end{itemize}

\end{document}
